\section{Huấn luyện theo phản hồi của con người}
\label{sec:ch02-huan-luyen-theo-phan-hoi}

Việc \tn{lời nhắc}{prompt} ảnh hưởng đến đầu ra không giải thích vì sao nhiều mô hình hội
thoại ưu tiên làm theo chỉ dẫn. Một quy trình phổ biến sau \tn{tiền huấn luyện}{pretraining} là
\tn{học tăng cường từ phản hồi của con người}{reinforcement learning from human feedback}.
Trong InstructGPT, quá trình này bắt đầu bằng fine-tune có giám sát trên lời
đáp mẫu, tiếp tục huấn luyện reward model từ các xếp hạng của người gán nhãn,
rồi tối ưu mô hình theo reward model bằng reinforcement learning
~\autocite{p04instructgpt}.

Hình~\ref{fig:ch02-huan-luyen-theo-phan-hoi} tách các dữ liệu và mô hình trong
quy trình đó. Sự tách biệt giúp tránh một suy luận quá mạnh: reward model ghi
lại một tín hiệu ưu tiên đã được học từ các xếp hạng, chứ không biến mọi câu
trả lời sau tối ưu thành đúng sự thật hoặc trung thực với một cơ chế bên trong.
Bài báo cũng ghi nhận các mô hình vẫn có thể không làm theo chỉ dẫn hoặc bịa
sự kiện~\autocite{p04instructgpt}.

\begin{figure}[htbp]
  \centering
  \begin{tikzpicture}[
  box/.style={draw=black!65, rounded corners=2pt, align=center,
    minimum width=4.4cm, minimum height=1.05cm, fill=black!5},
  data/.style={draw=black!65, rounded corners=2pt, align=center,
    minimum width=4.4cm, minimum height=0.9cm, fill=black!15},
  arrow/.style={-{Stealth[length=2mm]}, draw=black!65, thick}
]
  \node[box] (base) {Mô hình tiền huấn luyện};
  \node[data, right=11mm of base] (demos) {Lời đáp mẫu};
  \node[box, below=8mm of base] (sft) {Tinh chỉnh có giám sát};
  \node[data, right=11mm of sft] (rankings) {Xếp hạng của người gán nhãn};
  \node[box, below=8mm of sft] (reward) {Mô hình thưởng};
  \node[box, below=8mm of reward] (rl) {Tối ưu bằng học tăng cường};
  \node[box, below=8mm of rl] (chat) {Mô hình làm theo chỉ dẫn};

  \draw[arrow] (base.south) -- (sft.north);
  \draw[arrow] (demos.west) -- (sft.east);
  \draw[arrow] (rankings.west) -- (reward.east);
  \draw[arrow] (sft.south) -- (reward.north);
  \draw[arrow] (reward.south) -- (rl.north);
  \draw[arrow] (rl.south) -- (chat.north);
\end{tikzpicture}

  \caption{Sơ đồ khái niệm của quy trình InstructGPT. Lời đáp mẫu tạo dữ liệu
  fine-tune có giám sát; các xếp hạng của người gán nhãn tạo dữ liệu cho reward
  model; tín hiệu thưởng ấy được dùng trong bước tối ưu tiếp theo. Sơ đồ nêu
  đường đi của tín hiệu huấn luyện, không khẳng định rằng một điểm thưởng là
  lời giải thích cho từng token đầu ra~\autocite{p04instructgpt}.}
  \label{fig:ch02-huan-luyen-theo-phan-hoi}
\end{figure}

Đối với XAI, hệ quả là đối tượng cần kiểm tra có thể là mô hình sau toàn bộ
quy trình, không chỉ checkpoint tiền huấn luyện. Nếu một lời giải thích nói
rằng mô hình từ chối hay tuân theo một chỉ dẫn vì một đặc trưng của đầu vào,
ta cần nêu rõ checkpoint, mẫu lời nhắc và cách sinh đầu ra. Nếu không, hai
thí nghiệm có vẻ giống nhau có thể đang giải thích hai hành vi khác nhau.

\index{học tăng cường từ phản hồi của con người}%
